\par\Needspace{6\baselineskip}\medskip
\begin{center}\hypertarget{native-189-article141}{}\pdfbookmark[2]{141.}{article-141}141.\end{center}
Ex $P=\dfrac{n''}{n}$ atque aequatione III art. praec. sequitur
\[
(n+n'')\frac{P+a}{P+1}=bn'\frac{\sin(z-\sigma)}{\sin z}\mathpunct{\raisebox{.45ex}{;}}
\]
hinc vero et ex $Q=2\left(\dfrac{n+n''}{n'}-1\right)r'^3$ atque $r'=\dfrac{R'\sin\delta'}{\sin z}$ elicitur
\[
\sin z+\frac{Q\sin z^4}{2R'^3\sin\delta'^3}=b\frac{P+1}{P+a}\sin(z-\sigma)
\]
sive
\[
\begin{aligned}
\frac{Q\sin z^4}{2R'^3\sin\delta'^3}&=\left(b\frac{P+1}{P+a}-\cos\sigma\right)\sin(z-\sigma)\\[.5ex]
&\quad-\sin\sigma\cos(z-\sigma).
\end{aligned}
\]
Statuendo itaque brevitatis caussa
\[
\tag*{14)}\frac{1}{2R'^3\sin\delta'^3\sin\sigma}=c
\]
introducendoque angulum auxiliarem $\omega$ talem ut fiat
\[
\operatorname{tang}\omega=\frac{\sin\sigma}{b\dfrac{P+1}{P+a}-\cos\sigma}\mathpunct{\raisebox{.45ex}{,}}
\]
prodit aequatio
\[
\tag*{IV.}cQ\sin\omega\sin z^4=\sin(z-\omega-\sigma),
\]
\SourceLeaf{190}ex qua incognitam $z$ eruere oportebit. Ut angulus $\omega$ commodius computetur,
formulam praecedentem pro $\operatorname{tang}\omega$ ita exhibere conveniet
\[
\operatorname{tang}\omega=\frac{(P+a)\operatorname{tang}\sigma}{P\left(\dfrac{b}{\cos\sigma}-1\right)+\left(\dfrac{b}{\cos\sigma}-a\right)}\mathpunct{\cdot}
\]
Quamobrem statuendo
\[
\tag*{15)}\frac{\dfrac{b}{\cos\sigma}-a}{\dfrac{b}{\cos\sigma}-1}=d
\]
\[
\tag*{16)}\frac{\operatorname{tang}\sigma}{\dfrac{b}{\cos\sigma}-1}=e
\]
habebimus ad determinandum $\omega$ formulam simplicissimam
\[
\operatorname{tang}\omega=\frac{e(P+a)}{P+d}\mathpunct{\cdot}
\]
Computum quantitatum $a,b,c,d,e$ per formulas 11–16, a solis quantitatibus
datis pendentem, tamquam negotium quartum consideramus. Quantitates $b,c,e$
ipsae non erunt necessariae, verum soli ipsarum logarithmi.

Ceterum datur casus specialis, ubi haec praecepta aliqua mutatione indigent.
Quoties scilicet circulus maximus $BB''$ cum $A''B''$ coincidit, adeoque
puncta $B,B^\ast$ resp. cum $D',D$, quantitates $a,b$ valores infinitos nanciscerentur.
Statuendo in hoc casu
\[
\frac{R\sin\delta\sin(A'D''-\delta'+\sigma)}{R'\sin\delta'\sin(AD''-\delta)}=\pi,
\]
habebimus loco aequationis III hancce: $0=\pi n-\dfrac{n'\sin(z-\sigma)}{\sin z}\mathpunct{\raisebox{.45ex}{,}}$ unde faciendo
\[
\operatorname{tang}\omega=\frac{\pi\sin\sigma}{P+(1-\pi\cos\sigma)}
\]
eadem aequatio IV elicitur.

Perinde in casu speciali, ubi $\sigma=0$, fit $c$ infinita atque $\omega=0$, unde factor
$c\sin\omega$ in aequatione IV indeterminatus esse videtur: nihilominus revera
determinatus est, ipsiusque valor
\[
=\frac{P+a}{2R'^3\sin\delta'^3(b-1)(P+d)}\mathpunct{\raisebox{.45ex}{,}}
\]
uti levis attentio docebit. In hoc itaque casu fit
\[
\sin z=R'\sin\delta'\;\BareRoot[3]{\frac{2(b-1)(P+d)}{Q(P+a)}}\mathpunct{\cdot}
\]
